\section{Prompts and output scoring}
\label{app:prompts-scoring}

\paragraph{Arithmetic prompts.} The main few-shot prompt concatenates three
same-level ordinary-name demonstration blocks, separated by blank lines, followed
by the test equations and query. Each demonstration contains its equations and
query on one line and its answer or calculation on the next. Answer-only output
has the form \texttt{pen=6}. Calculation output restates equations, substitutes
resolved values and reports each evaluated value in dependency order, ending
with the queried variable. The test prompt ends immediately after its input line
and a newline. Demonstrations reserve a name pool disjoint from test names;
letter-name replication uses letter-name demonstrations. Task levels and
name conditions are in Appendix~\ref{a-app:task}; additional equation and final-query
formats are in Appendix~\ref{a-app:round7-controls}.

\paragraph{Code prompts.} Both main code formats concatenate three ordinary-name
worked examples, separated by blank lines, followed by the test function, its
call and \texttt{Trace:}. There is no additional instruction or interpreter
execution. Each example contains the program and input call, a \texttt{Trace:}
line and an \texttt{Answer:} line. Table~\ref{tab:code-demonstrations} gives the
complete level-5 demonstration set used for Figure~\ref{fig:example}. Assemble
one format by choosing its trace for each example. Both formats then append the
same test program from that figure and the same \texttt{Trace:} prefix. Traces
wrap here for typesetting; each is one line in the prompt. The additional formats
are specified in Appendix~\ref{c-app:formats}.

\begin{table*}[t]
\centering\small
\begin{tabular}{@{}p{.32\textwidth}p{.25\textwidth}p{.37\textwidth}@{}}
\toprule
Program and input & Values-only trace & Expression-and-value trace \\
\midrule
\texttt{def f(xs):}\newline \hspace*{1em}\texttt{u = min(xs)}\newline \hspace*{1em}\texttt{t = u + 1}\newline \hspace*{1em}\texttt{k = t * 2}\newline \hspace*{1em}\texttt{return k}\newline \texttt{f([0,\allowbreak{} 6,\allowbreak{} 5])} & \texttt{u = 0,\allowbreak{} t = 1,\allowbreak{} k = 2}\newline \texttt{Answer: 2} & \texttt{u = min(xs) = 0,\allowbreak{} t = u + 1 = 1,\allowbreak{} k = t * 2 = 2}\newline \texttt{Answer: 2} \\
\addlinespace
\texttt{def f(xs):}\newline \hspace*{1em}\texttt{z = len(xs)}\newline \hspace*{1em}\texttt{r = z * 2}\newline \hspace*{1em}\texttt{k = r - 1}\newline \hspace*{1em}\texttt{return k}\newline \texttt{f([2,\allowbreak{} 4,\allowbreak{} 1])} & \texttt{z = 3,\allowbreak{} r = 6,\allowbreak{} k = 5}\newline \texttt{Answer: 5} & \texttt{z = len(xs) = 3,\allowbreak{} r = z * 2 = 6,\allowbreak{} k = r - 1 = 5}\newline \texttt{Answer: 5} \\
\addlinespace
\texttt{def f(xs):}\newline \hspace*{1em}\texttt{p = len(xs)}\newline \hspace*{1em}\texttt{ww = p + 2}\newline \hspace*{1em}\texttt{u = ww // 2}\newline \hspace*{1em}\texttt{return u}\newline \texttt{f([3,\allowbreak{} 6,\allowbreak{} 0])} & \texttt{p = 3,\allowbreak{} ww = 5,\allowbreak{} u = 2}\newline \texttt{Answer: 2} & \texttt{p = len(xs) = 3,\allowbreak{} ww = p + 2 = 5,\allowbreak{} u = ww // 2 = 2}\newline \texttt{Answer: 2} \\
\addlinespace
\bottomrule
\end{tabular}
\caption{Complete code demonstration set 7; prepend \texttt{Trace:} to each trace.}
\label{tab:code-demonstrations}
\end{table*}


\paragraph{Parsing and denominators.} For code few-shot output, parsing stops at
the first blank line, so an invented subsequent problem cannot supply the answer.
The final answer is the last \texttt{Answer: <integer>} in that block; answer-only
continuations may instead start with the integer directly. A trace value is the
last integer immediately following an equals sign in the target variable's
first assignment clause, before \texttt{Answer:}. This retains the evaluated
value in an expression-and-value trace rather than an operand. A comma followed
by another assignment ends that clause. Interactive transcripts read the integer
following the queried variable; comments read its end-of-line value. Free prose
uses the final \texttt{Answer:} occurrence, or the last integer on the final
nonempty line when that marker is absent.

For arithmetic few-shot output, the answer is the last standalone
\texttt{<queried-name>=<integer>} clause before the first blank line. Free-text
arithmetic uses a queried-name assignment or a final integer fallback. The
reported main arithmetic name-error rate grades this final answer. Code
intermediate correctness grades the first assignment against the program's
executed value; final correctness grades the return value. Parse failures count
as incorrect for accuracy and do not match a name-suggested digit. They remain
in every behavioral denominator. Error-conditioned probes separately use the
observed wrong-write cohort with verified generated prefixes.

\paragraph{Sampling and uncertainty.} Main behavior uses demonstration seeds
7, 11 and 13, with three examples per set and greedy full-vocabulary decoding.
The arithmetic panel contains \NUM{a-n-models} checkpoints and
\NUM{a-n-test-sets} matched sets per task level. The code panel contains
\NUM{c-n-models} instruction-tuned checkpoints, with an additional 34B check,
and 1,000 matched problems per target at levels 3 and 5. The central code cohort
has \NUM{c-probe-n-test} unique length-to-sum computations per set. Checkpoints
and prompts differ between tasks; their results are not pooled.

Resampling retains both names and all demonstration repeats of a computation.
Behavioral reliability requires a consistent sign in all three sets and a 95\%
interval excluding zero. The small-effect criterion uses a 90\% interval within
the prespecified two-point margin. Code timing and matched correctness intervals
use 4,000 draws; the code equivalence audit uses 2,000. Intervals are pointwise,
and probe intervals condition on the fixed independent layer selectors. These
procedures are detailed in Appendices~\ref{c-app:stats},
\ref{c-app:prompt-comparison} and \ref{app:format-correctness}; arithmetic repeated-
computation checks are in Appendix~\ref{a-app:round7-paired}.

\paragraph{Probe fitting and records.} Predictors use three initialization seeds,
0, 1 and 2, and full-batch SGD at learning rate $10^{-3}$ for 10,000 epochs.
Training and evaluation are disjoint in computation, ignoring identifier
spelling. Five-fold layer selection excludes the held-out computation. Cached
prefix checks exclude the future answer token. Optimizer checks, controls,
software versions and data records are in Appendices~\ref{a-app:recipe},
\ref{c-app:round6-probes} and \ref{c-app:reproducibility}.
